\documentclass[10pt,a4paper]{amsart}
\usepackage[margin=19mm]{geometry}
\usepackage{amsmath,amssymb,mathtools}
\usepackage[scr=rsfs,cal=euler]{mathalfa}
\usepackage{xfrac}
\usepackage[unicode,hidelinks,bookmarks=false]{hyperref}
\newcommand{\ie}{i.e.\ }
\newcommand{\cf}{cf.\ }
\newcommand{\resp}{respectively\ }
\newcommand{\Subsection}{\subsection}
\providecommand{\nobreakdash}{\nobreak\hbox{-}}
\setlength{\emergencystretch}{2em}
\multlinegap0pt
\usepackage{amssymb}
\newtheorem{theorem}{Theorem}
\title{Natural Orderings of Triangle Centers}
\author{Stanley Rabinowitz}
\date{Source: arXiv:2603.15698v1 (CC BY 4.0)}
\begin{document}
\maketitle

\noindent\emph{Source and licence.} Natural Orderings of Triangle Centers. Version arXiv:2603.15698v1 (CC BY 4.0), distributed by arXiv under the Creative Commons Attribution 4.0 International licence, http://creativecommons.org/licenses/by/4.0/, as recorded in the arXiv licence metadata for this exact version and shown on the abstract page https://arxiv.org/abs/2603.15698v1. Full licence notice: \texttt{licenses/arxiv-2603.15698-CC-BY-4.0.txt}.\par\smallskip

\noindent\textit{Editorial scope.} Counted unit: the article's theorem theorem (source0.tex lines 208--212). The article's complete original proof of it is delivered with it (source0.tex lines 214--226, one \texttt{proof} environment of 13 lines). No mathematics is written, repaired or re-derived: the statement and the proof are the article's own lines, in the article's own notation, and no retained line is substituted or re-flowed.  Retained uncounted as the source-local context the proof needs: lines 197--202.  One declared adaptation: exactly 1 ancillary strings inside the retained spans are omitted (image-inclusion commands and, where the source repeats a label, the repeated label definitions) and no other line differs from the source; the figure environments, with their placement, captions and labels, are kept, and the package records the omitted strings in fidelity receipts bound to the affected bodies.  No retained line contains a citation, so the wrapper carries no bibliography and no bibliography entry.  The retained lines contain the article's own diagram code, which the wrapper typesets with the standard tikz packages; no image file is delivered and the package declares no asset dependency.  Editorial formatting only, in addition: an AMS \texttt{amsart} preamble that restates the article's own declarations and macros used by the retained lines, namely the environments \texttt{theorem} on separate counters, together with the standard packages those lines need.  The retained environments print in this document as Theorem 1 in order of appearance, whereas the article prints the same items with the numbering of its own full document.  The complete native source of the article as distributed in the arXiv e-print for this exact version is bundled unchanged as \texttt{sources/196497/source0.tex}.
\begin{figure}[h!t]
\centering

\caption{Seven regions formed by the sidelines of a triangle}
\label{fig:7areas}
\end{figure}

\begin{theorem}
Let $P$ have barycentric coordinates $(u:v:w)$ with respect to $\triangle ABC$.
Then $P$ is inside angle $A$ if and only if
$$v(u+v+w)>0\qquad\textrm{and}\qquad w(u+v+w)>0.$$
\end{theorem}

\begin{proof}
Since barycentric coordinates are homogeneous, the coordinates $(u:v:w)$ and $(uk:vk:wk)$
represent the same point for any $k\neq 0$.
The normalized barycentric coordinates for $P$ are therefore
$$\left(\frac{u}{u+v+w}:\frac{v}{u+v+w}:\frac{w}{u+v+w}\right).$$
Note that $u+v+w\neq 0$ since $P$ is not on the line at infinity.
From Figure~\ref{fig:7areas}, we see that for normalized coefficients, $P=(p:q:r)$ is inside angle $A$
if and only if $q>0$ and $r>0$.
In terms of $u$, $v$, and $w$, this means that
$v/(u+v+w)>0$ and $w/(u+v+w)>0$.
Multiplying each inequality by the positive quantity $(u+v+w)^2$ preserves
the sense of the inequality and gives us the desired result.
\end{proof}

\end{document}
